\documentclass{article}
	\usepackage[UTF8]{ctex}
	\usepackage{amsmath}
	\usepackage{graphicx}
	\usepackage{underscore}
	\usepackage{geometry}
	\usepackage{anysize}
	\date{}
	\title{全新比赛机制 \\ [2ex] \begin{large} 
	时间限制：C/C++ 1秒，其他语言2秒 \\
	空间限制：C/C++ 256.000MB，其他语言512.000MB 	\\
	64bit IO Format: \%lld
 			\end{large} }
	\geometry{top=10mm,bottom=15mm,left=25mm,right=25mm}
	
	\begin{document}
		\maketitle \vspace{-50pt}
		
		\section*{题目描述}
		\indent LZU_ACM小队成功地打进2018年ACM/ICPC(ACM International Collegiate Programming Contest)的现场赛啦！正当他们满怀着激动地心情到达目的地时，却被告知比赛规则变了。原来一个队伍只能使用一台电脑，而现在可以使用三台电脑，同时队员之间不得相互配合。幸运地是，由于他们做题经验丰富，并且掌握一些不为人知的玄学算法，他们已经提前知道了题目数量为13题，以及不同同学完成不同题目所需的时间。他们想知道完成所有题目最快的时间。\\
		\indent 针对这个问题，LZU_ACM小队的成员各有不同的求解方法，但是他们都被一些事情(与AK大佬合影)耽误了，于是他们把这个问题交给你。

		\section*{输入描述}
			\noindent 输入13*3的矩阵。第$i$($1 \le i \le 13$)行输入三个整数$t_{i1},t_{i2},t_{i3}$($1 \le t_{ij}\le 20$,$1 \le j \le 3$)，代表第j个同学完成第i道题的时间为$t_{ij}$。所有数据采用随机生成，随机方式为 $rand()\%20+1$。

		\section*{输出描述}
			\noindent 输出一个整数，为完成所有题目最快的时间。
			
		\section*{示例1}
			\subsection*{输入}
\noindent
1 2 3	\\
1 2 3	\\
1 2 3	\\
1 2 3	\\
3 1 2	\\
3 1 2	\\
3 1 2	\\
3 1 2	\\
2 3 1	\\
2 3 1	\\
2 3 1	\\
2 3 1	\\
1 3 5

			\subsection*{输出}
\noindent
5

			\subsection*{说明}
			\noindent 第一个同学完成第1,2,3,4,13道题，第二个同学完成第5,6,7,8道题，第三个完成第9,10,11,12道题。
			
		\section*{备注}

	\end
	{document}